% Pista alg-25i-q2 · Àlgebra
% Procedència: PAU Catalunya 2025 · Convocatòria de juny (incidències) · Sèrie 4 · Exercici 2
\documentclass[11pt,a4paper]{article}
\usepackage{pista}
\pistainfo{Catalunya 2025}{Convocatòria de juny (incidències)}{Sèrie 4}

\begin{document}

\section*{\textcolor{blauPista}{Exercici 2}}

\noindent Considereu el sistema d'equacions
\[
\left.\begin{array}{r}
x + 3y + az = 1 \\
-2x - 6y + 2z = 2a \\
ax + 3y + z = 3
\end{array}\right\},
\]
on $a$ és un paràmetre real, i siguin $\pi_{1}$ el pla determinat per la primera equació, $\pi_{2}$ el determinat per la segona, i $\pi_{3}$ el determinat per la tercera.

\subsection*{\textit{a)} \normalfont Discutiu el sistema en funció del valor de $a$. \hspace*{\fill}\textnormal{[1~punt]}}

\pista{Escriu la matriu $A$ i la matriu ampliada $A^{'}$, i calcula el determinant de la matriu $A$, aplicant la regla de Sarrus. A continuació, estudia el sistema d'equacions. D'una banda, has d'estudiar a part cada valor del paràmetre $a$ que anul·la el determinant de la matriu $A$. D'altra banda, la resta de valors que no anul·len el determinant fan que el sistema sigui compatible determinat, és a dir, que té solució única.

\medskip
Per estudiar cada valor del paràmetre $a$ que anul·la el determinant, esglaona la matriu ampliada, aplicant el mètode de Gauss. Et convé començar amb $F_{2} \to F_{2} + 2F_{1}$, veus per quin motiu?}

\subsection*{\textit{b)} \normalfont Descriviu la posició relativa dels tres plans $\pi_{1}$, $\pi_{2}$ i $\pi_{3}$ en funció del valor de $a$. \hspace*{\fill}\textnormal{[0,75~punts]}}

\pista{Per a cada cas de la discussió de l'apartat \textit{a)}, compara els plans de dos en dos. Si els coeficients de $x$, $y$ i $z$ de dos plans són proporcionals, aquests dos plans són paral·lels. Si, a més, els termes independents també ho són, llavors els dos plans són coincidents.

\medskip
Llavors, per a cada valor del paràmetre $a$, digues si els tres plans es tallen en un punt, en una recta, o si no tenen cap punt en comú, i quins plans són paral·lels o coincidents.}

\subsection*{\textit{c)} \normalfont En algun cas la intersecció dels tres plans és una recta? En cas afirmatiu, digueu per a quin valor del paràmetre, i doneu un vector director i un punt d'aquesta recta. \hspace*{\fill}\textnormal{[0,75~punts]}}

\pista{Mira l'apartat \textit{b)}: per a quin valor del paràmetre $a$ els tres plans es tallen en una recta? Per a aquest valor, aprofita la matriu esglaonada que has obtingut a l'apartat \textit{a)}. A continuació, escriu simplement \mbox{$z = t$} (és un paràmetre lliure), i finalment aïlla $x$ i $y$ en funció de $t$.

\medskip
Escriu la solució separant la part que no depèn de $t$:
\[
(x, y, z) = (\ldots) + t\,(\ldots)
\]
La primera part és un punt de la recta, i la segona, un vector director.}

\end{document}
